% SPDX-License-Identifier: GPL-3.0-or-later
% Copyright (C) 2026 Alessandro Greco
\chapter{Наблюдение за трафиком: Wireshark и журналы}\label{cap:20-traffico}
\aggiornatoal{23 сентября 2026 г.}{tor 0.4.9.13; tshark (Debian trixie); лабораторная сеть}

\begin{obiettivi}
К концу этой главы вы сможете:
\begin{itemize}[nosep]
  \item перехватывать трафик Tor с помощью \cmd{tcpdump} и анализировать его через \cmd{tshark} или Wireshark;
  \item отличать то, что видит внешний наблюдатель, от того, что остаётся скрытым;
  \item расшифровывать TLS на \emph{собственном} канале с помощью сессионных ключей и исследовать его ячейки (Лабораторная работа~12);
  \item пошагово прослеживать протокол по журналам tor.
\end{itemize}
\end{obiettivi}

В предыдущих главах мы исследовали протокол \enquote{изнутри}, создавая код, который разговаривает на нём. Здесь мы наблюдаем его \enquote{снаружи}, как это делал бы тот, кто следит за сетью. У такого упражнения две противоположные цели: понять, что атакующий может вывести из трафика, и получить независимый способ проверки собственной реализации.

\begin{attenzione}[Только собственный трафик и только в лабораторной среде]
Перехват чужого трафика без разрешения незаконен в Италии, как и в большинстве стран (\cref{cap:25-legale}), и противоречит этике исследований анонимности (\cref{sec:rp-perche}). Все перехваты в этой главе относятся либо к приватной сети книги, либо к соединениям, которые мы открываем сами.
\end{attenzione}

% ---------------------------------------------------------------------------
\section{Что может видеть наблюдатель}\label{sec:traffico-osservatore}

Тот, кто наблюдает соединение между клиентом и его входным реле (guard), например интернет-провайдер, может видеть:
\begin{itemize}
  \item IP-адреса и порты: клиент подключается к адресу, который, если это не мост (bridge), присутствует в публичном консенсусе и поэтому распознаётся как реле Tor;
  \item рукопожатие TLS: сертификаты передаются в зашифрованном виде при TLS~1.3, но некоторые свойства начального сообщения клиента остаются видимыми;
  \item размеры и временные характеристики записи TLS.
\end{itemize}
Наблюдатель не видит содержимое ячеек, цепи или адреса назначения. Однако именно размеры и временные характеристики лежат в основе атак \emph{website fingerprinting} и корреляции (\cref{cap:26-attacchi}). Зашифрованное содержимое не означает, что трафик вообще ничего не раскрывает.

В \cref{sec:lab-censura} мы сравнивали обычный Tor-трафик, распознаваемый как TLS, с трафиком obfs4, не имеющим узнаваемой структуры. Здесь мы пойдём ещё дальше и заглянем внутрь TLS.

% ---------------------------------------------------------------------------
\section{Инструменты}\label{sec:traffico-strumenti}

\begin{description}
  \item[\cmd{tcpdump}] Перехватывает пакеты на сетевом интерфейсе и сохраняет их в файл \texttt{pcap}. Например, внутри лабораторного контейнера можно перехватить трафик к ORPort узла \texttt{test004r}:
\begin{lstlisting}[language=bash,numbers=none]
$ docker exec -u 0 torbook-network tcpdump -i lo \
    -w /lab/net/tls.pcap "tcp port 5104"
\end{lstlisting}
  \item[Wireshark и \cmd{tshark}] Анализируют захваченные пакеты: первый через графический интерфейс, второй из командной строки. Они распознают TLS, но официального анализатора протокола Tor нет. Даже после расшифрования ячейки Tor выглядят как необработанные байты. В роли анализатора будет выступать наш модуль \file{cells.py}.
\end{description}

\begin{figure}[htbp]
\centering
% SPDX-License-Identifier: GPL-3.0-or-later
% Copyright (C) 2026 Alessandro Greco
%% Альтернативное описание для EPUB.
\altfig{Упрощённая схема Tor поверх TCP и TLS. Внешний наблюдатель видит IP-адреса, порты, объём трафика, временные характеристики и длины TLS-записей, но без сессионных ключей TLS не видит ячейки. В одной TLS-записи могут находиться несколько ячеек, а одна ячейка может быть разделена между записями; их границы не обязаны совпадать. При протоколе канала версии 4 и выше ячейка фиксированной длины составляет 514 байт: заголовок 5 байт и тело 509 байт; существуют также ячейки переменной длины. Послойное шифрование ключами цепи применяется к телу ячеек RELAY, а не ко всем типам, например PADDING и VPADDING. Ключи TLS позволяют увидеть структуру ячеек, но не расшифровывают автоматически тело RELAY.}
% Схема инкапсуляции: границы TLS records и Tor cells не обязаны совпадать.
\begin{tikzpicture}[font=\sffamily\scriptsize,
  liv/.style={draw,rounded corners=2pt,minimum height=0.85cm,align=center},
  cel/.style={draw,fill=torviola!12,minimum height=0.7cm,minimum width=1.5cm,font=\sffamily\tiny}]
  \node[liv,fill=torgrigio!10,minimum width=10.6cm] (tcp) at (0,2.6) {TCP: сегменты переменной длины};
  \node[liv,fill=torverde!12,minimum width=9.4cm] (tls) at (0,1.35) {TLS: границы не равны границам cells};
  \foreach \i/\x in {1/-3.6,2/-1.8,3/0,4/1.8,5/3.6} \node[cel] (c\i) at (\x,0) {514\,B};
  \node[anchor=west,font=\sffamily\tiny,text=torgrigio] at (4.6,0) {фикс. v4+};
  \node[liv,fill=avviso!10,minimum width=4.6cm] (rel) at (-1.7,-1.4) {RELAY: тело 509\,Б\\послойно зашифровано};
  \node[liv,fill=torgrigio!8,minimum width=4.6cm] (pad) at (3.3,-1.4) {PADDING / VPADDING\\не RELAY};
  \draw[torgrigio] (c1.south west) -- (rel.north west); \draw[torgrigio] (c3.south east) -- (rel.north east);
  \node[anchor=west,align=left,font=\sffamily\tiny,text=torgrigio] at (5.8,2.6) {IP, порты, объём\\и время видимы};
  \node[anchor=west,align=left,font=\sffamily\tiny,text=torgrigio] at (5.8,1.35) {длина TLS-записей\\видна на канале};
  \node[anchor=west,align=left,font=\sffamily\tiny,text=torgrigio] at (5.8,0) {ячейки видны после\\расшифрования TLS};
  \node[anchor=west,align=left,font=\sffamily\tiny,text=torgrigio] at (5.8,-1.4) {RELAY: нужны\\ключи цепи};
  \node[font=\sffamily\tiny,text=torgrigio,align=center] at (0,-2.55) {Схема условна: границы TLS-записей и ячеек могут не совпадать.};
\end{tikzpicture}

\caption{Схема канала Tor: сегменты TCP передают записи TLS, границы которых не обязаны совпадать с границами ячеек. Наблюдатель видит IP-адреса, порты, длины и время; после расшифрования TLS можно выделить ячейки. При версии протокола канала v4+ фиксированная ячейка занимает 514 байт (заголовок 5, тело 509), но существуют также ячейки переменной длины. Послойно шифруются тела ячеек \texttt{RELAY}, а ячейки заполнения не являются сообщениями RELAY.}
\label{fig:incapsulamento}
\end{figure}

\subsection{Расшифрование собственного TLS}

TLS использует сессионные ключи, которыми внешний наблюдатель не располагает. Однако сторона, \emph{участвующая} в соединении, может экспортировать эти ключи. Браузеры и приложения Python могут записывать ключи сессии в файл, если включено их журналирование. Наш модуль \file{channel.py} явно устанавливает \texttt{SSLContext.keylog\_filename}, когда определена переменная \texttt{SSLKEYLOGFILE}. Wireshark и \cmd{tshark} используют этот файл для расшифрования перехвата. Скрипт \file{channel.py} из Лабораторной работы~1 следует этому соглашению: установите переменную перед запуском. Файл содержит строки вида:
\begin{lstlisting}[numbers=none]
# TLS secrets log file, generated by OpenSSL / Python
SERVER_HANDSHAKE_TRAFFIC_SECRET 4878224d3f9fb9aedb96d... 0da89dff6feb...
\end{lstlisting}
Они содержат случайное клиентское значение и секреты каждой фазы TLS~1.3 для соответствующей сессии.

\begin{approfondimento}[Почему это не атака на Tor]
Расшифрование TLS на собственном канале не раскрывает ничего, чего клиент уже не знал. Даже при наличии сессионных TLS-ключей ячейки \cella{RELAY} остаются послойно зашифрованными ключами отдельных переходов (\cref{cap:10-cifratura-relay}). Для их чтения потребовались бы ключи цепи, которыми располагают только клиент и соответствующие реле. TLS — внешняя оболочка многослойного шифрования, а не единственный уровень защиты.
\end{approfondimento}

% ---------------------------------------------------------------------------
\section{Лабораторная работа: ячейки внутри TLS}\label{sec:lab-traffico}

\begin{laboratorio}[12]{Перехват, расшифрование и чтение канала}
\textbf{Цель.} Перехватить канал, открытый нашим клиентом к реле, расшифровать его с помощью сессионных ключей и восстановить последовательность ячеек через \file{labs/ru/код/torlab/traffic.py}.\\
\textbf{Предварительные условия.} Лабораторные работы~1--3; лабораторная сеть должна быть запущена; в контейнере должны быть установлены \cmd{tcpdump} и \cmd{tshark}.\\
\textbf{Оценочное время.} 1--2 часа.
\end{laboratorio}

Для лабораторной работы~12 запустите приватную сеть и откройте два терминала из корня репозитория. Следующую команду захвата выполняйте \emph{на хосте}: она запускает \cmd{tcpdump} от root \emph{внутри лабораторного контейнера}, поскольку обычная оболочка Makefile работает с правами непривилегированного пользователя. Пока во втором терминале работает клиент, оставьте захват активным, затем остановите его сочетанием Ctrl+C. Файл сохраняется в общем томе \file{labs/ru/лаборатория/net/}.

\paragraph{Шаг 1: перехват.} Запустите перехват внутри контейнера. Затем с хоста откройте канал к \texttt{test004r} и постройте одноузловую цепь, одновременно экспортируя сессионные ключи:
\begin{lstlisting}[language=bash,numbers=none]
$ PYTHONPATH=labs/ru/код SSLKEYLOGFILE=labs/ru/лаборатория/net/tls-keys.log python3 -c "
from torlab.laboratory import relay
from torlab.circuit import build
build([relay('004r')], cc=False)"
\end{lstlisting}

Команду Python выполняйте из \emph{корня репозитория} в Linux с сетевым режимом хоста. \texttt{PYTHONPATH} делает пакет \file{torlab} из \file{labs/ru/код/} доступным для импорта; \file{labs/ru/код/torlab/channel.py} явно считывает \texttt{SSLKEYLOGFILE} и включает экспорт секретов TLS. Файл содержит секретные ключевые данные: не публикуйте его и удалите после завершения опыта.

\paragraph{Шаг 2: определение нашего соединения.} Другие лабораторные реле тоже подключаются к тому же порту, но их сессионные ключи отличаются. Попросим \cmd{tshark} проследить каждый поток TCP после попытки расшифрования и сохраним тот, который успешно расшифровывается:
\begin{lstlisting}[language=bash,numbers=none]
$ docker exec torbook-network tshark -r /lab/net/tls.pcap \
    -o tls.keylog_file:/lab/net/tls-keys.log \
    -q -z follow,tls,raw,2
\end{lstlisting}
Номер потока TCP \texttt{2} относится лишь к предыдущей демонстрационной записи. Найдите возможные номера потоков с помощью \cmd{tshark} и \texttt{-T fields -e tcp.stream}, затем выберите расшифровываемый поток. Приведённый ниже вывод исторический: размеры сертификатов и идентификаторы цепей меняются при новом запуске.
Вывод содержит расшифрованные данные обоих направлений в шестнадцатеричном виде; одно из направлений отображается с отступом табуляцией.

\paragraph{Шаг 3: исследование ячеек.} Функция \texttt{cells\_from\_follow} восстанавливает оба потока и передаёт их в \texttt{decode} из Лабораторной работы~2. Каждое направление имеет собственное согласование версии: первой ячейкой является \cella{VERSIONS} с двухбайтовым CircID. Четырёхбайтовые CircID используются только после согласования. Результат:
\begin{lstlisting}[numbers=none]
клиент -> реле  CircID=0x0000 VERSIONS тело=4 байт
реле   -> клиент CircID=0x0000 VERSIONS тело=6 байт
реле   -> клиент CircID=0x00000000 CERTS тело=1443 байт
реле   -> клиент CircID=0x00000000 AUTH_CHALLENGE тело=36 байт
реле   -> клиент CircID=0x00000000 NETINFO тело=509 байт
клиент -> реле  CircID=0x00000000 NETINFO тело=509 байт
клиент -> реле  CircID=0xc244f1f6 CREATE2 тело=509 байт
реле   -> клиент CircID=0xc244f1f6 CREATED2 тело=509 байт
\end{lstlisting}
При внешнем наблюдении это в точности последовательность из \cref{fig:sessione} и \cref{cap:07-canali}. Ячейка клиента \cella{VERSIONS} имеет тело длиной четыре байта, поскольку \file{channel.py} предлагает только версии 4 и 5. У Circuit ID \texttt{0xc244f1f6} старший бит равен 1, поскольку канал открыл клиент (\cref{sec:circid}).

\begin{attenzione}[Две ошибки в нашей первой версии]
Первоначальный анализ содержал две ошибки, которые обнаружились при сравнении вывода с ожидаемой последовательностью. Мы использовали одну согласованную версию протокола для обоих направлений и поэтому неправильно интерпретировали ячейку \cella{VERSIONS} от реле. Кроме того, мы назначили блоки неверным направлениям, потому что предположили, какое именно направление \cmd{tshark} выводит с отступом. Правильное правило простое: первый блок всегда является клиентской ячейкой \cella{VERSIONS}.
\end{attenzione}

% ---------------------------------------------------------------------------
\section{Журналы tor}\label{sec:traffico-log}

Перехват пакетов показывает, \emph{что} передаётся по сети. Журналы tor объясняют, \emph{почему}. На уровне \texttt{info} tor записывает этапы построения цепей, выбор входных реле, загрузку и публикацию дескрипторов. На уровне \texttt{debug} записи доходят до отдельных ячеек. Мы уже использовали журналы, чтобы проверить, к какому HSDir обращался клиент (\cref{sec:lab-onion}) и когда onion-сервис ротировал свои дескрипторы.

Несколько практических рекомендаций:
\begin{itemize}
  \item Уровень \texttt{debug} создаёт огромные объёмы текста. Включайте его только для конкретного домена, например \texttt{Log [circ]debug file circ.log}, чтобы ограничить подробные сообщения событиями цепей.
  \item В реальной сети \opt{SafeLogging} маскирует адреса. Журналы клиента никогда не должны содержать посещаемые им назначения.
  \item Журналы публичного реле могут содержать персональные данные других людей с юридическими последствиями, рассматриваемыми в \cref{cap:25-legale}; сохраняйте только строго необходимую информацию.
\end{itemize}

% ---------------------------------------------------------------------------
\section*{Итоги}
\begin{itemize}
  \item Наблюдатель соединения видит адреса, рукопожатие TLS, размеры пакетов и временные характеристики, но не ячейки Tor и не адреса назначения.
  \item \texttt{SSLKEYLOGFILE} позволяет расшифровать TLS на собственном канале и исследовать его ячейки; ячейки \cella{RELAY} при этом всё равно остаются послойно зашифрованными.
  \item Анализ перехваченного трафика предоставляет независимый способ проверить реализацию.
  \item Журналы tor уровней \texttt{info} и \texttt{debug} объясняют решения программы; при реальном использовании объём логирования следует минимизировать и защищать.
\end{itemize}

% ---------------------------------------------------------------------------
\section*{Упражнения}

\begin{esercizio}[1]\label{es:20-dimensioni}
В перехвате загрузки через Tor не все записи TLS имеют длину, кратную 514 байтам. Почему, если ячейки имеют фиксированный размер?
\end{esercizio}

\begin{esercizio}[2]\label{es:20-relay}
Расширьте Лабораторную работу~12, расшифровав ячейки \cella{RELAY} собственной цепи с помощью ключей из \file{circuit.py}. Какие сообщения реле видны при построении трёхузловой цепи?
\end{esercizio}

\begin{esercizio}[2]\label{es:20-padding}
Перехватывайте канал лабораторного клиента в течение одной минуты без создания прикладного трафика. Подсчитайте записи TLS в каждом направлении и сопоставьте наблюдения с заполнением канала (channel padding) из \cref{cap:07-canali}.
\end{esercizio}

\begin{esercizio}[3]\label{es:20-dissector}
Напишите небольшой анализатор Wireshark на Lua, который распознаёт ячейки Tor в расшифрованных данных и отображает их CircID и команду.
\end{esercizio}

% ---------------------------------------------------------------------------
\section*{Дополнительная литература}
\begin{itemize}
  \item Документация Wireshark по расшифрованию TLS с использованием файлов журнала сессионных ключей.
  \item Раздел \opt{Log} руководства tor, описывающий домены журналирования и уровни детализации.
\end{itemize}
